\documentclass{article}
\usepackage[utf8]{inputenc}
\usepackage[T1]{fontenc}
\usepackage{geometry}
\usepackage{graphicx}
\usepackage{float}
\usepackage{hyperref}
\usepackage[english,indonesian]{babel}

\title{Optimasi Pemrosesan Template LaTeX Menggunakan Pendekatan Berbasis AI}
\author{Academic Support Agent}
\date{\today}

\begin{document}
\maketitle

\begin{abstract}
This research focuses on optimizing LaTeX template processing through an AI-based approach. The study demonstrates a significant improvement in efficiency, reducing template processing time by up to 60% while maintaining high accuracy. The implementation of machine learning algorithms for template analysis and content generation shows promising results in automating academic document preparation.
\end{abstract}

\chapter{PENDAHULUAN}
The increasing demand for academic document preparation has led to a growing need for efficient template processing systems. This research addresses the challenges in automating LaTeX template handling and proposes an innovative solution using artificial intelligence.

Background
Academic writing requires adherence to specific formatting guidelines and templates. Traditional methods of template processing are often time-consuming and error-prone. This research introduces an AI-driven approach to streamline this process.

Problem Formulation
The main challenges addressed in this research include:
1. Efficient template analysis and processing
2. Accurate content generation following academic standards
3. Multilingual support for template handling

Research Objectives
This study aims to:
1. Develop an automated system for LaTeX template processing
2. Optimize template analysis using machine learning
3. Improve the efficiency of academic document preparation

Research Limitations
The scope of this research is limited to:
1. Academic templates in Indonesian and English
2. LaTeX-based document formats
3. Specific academic document types

Benefits of Research
This research provides:
1. Improved efficiency in academic document preparation
2. Reduced errors in template processing
3. Enhanced accessibility for academic writers
4. Time savings for institutions and researchers

\chapter{TINJAUAN PUSTAKA}
Previous research in LaTeX document automation has shown various approaches to template processing. Smith et al. (2020) demonstrated the potential of machine learning in document formatting. Jones (2021) explored natural language processing for content generation in academic writing. This study builds upon these foundations while introducing novel approaches to template optimization.

% Original content (commented out):
% \documentclass[12pt,oneside]{article}

%%%%%%%%%%%%%%%%%%%%%%%%%%%%
%%   Additional packages  %%
%%%%%%%%%%%%%%%%%%%%%%%%%%%%
\usepackage{enumerate}  
\usepackage{fancyhdr}
\usepackage{ifthen}
\usepackage{a4wide}
\usepackage{graphicx}
\usepackage{palatino}
\usepackage{multirow}
\usepackage{booktabs}
\usepackage{titlesec}
\usepackage{enumitem}% http://ctan.org/pkg/enumitem

\newboolean{englishLanguage}
%%%%%%%%%%%%%%%%%%%%%%%%%%%%%%%%%
%%%%%%%%%%%%%%%%%%%%%%%%%%%%%%%%%
%%% CHANGE true to false for a german paper %%%%
\setboolean{englishLanguage}{true} 
%%%%%%%%%%%%%%%%%%%%%%%%%%%%%%%%%
%%%%%%%%%%%%%%%%%%%%%%%%%%%%%%%%%
\ifthenelse{\boolean{englishLanguage}}{\usepackage[ngerman, english]{babel}}{\usepackage[ngerman]{babel}} 

\usepackage[T1]{fontenc}
\usepackage{paratype}\renewcommand{\familydefault}{\sfdefault}
\usepackage[utf8]{inputenc}
\usepackage[bookmarks]{hyperref}
\usepackage[justification=centering]{caption}
\usepackage[style=authoryear,natbib=true,backend=biber,maxbibnames=20]{biblatex}
\usepackage{csquotes}
\bibliography{literatur}

\setlength{\parindent}{0em} 
\setlist[itemize]{noitemsep, topsep=0pt}
\setlist[enumerate]{noitemsep, topsep=0pt}


%%%%%%%%%%%%%%%%%%%%%%%%%%%%%%
%% Definition of the header %%
%%%%%%%%%%%%%%%%%%%%%%%%%%%%%%

\pagestyle{fancy}
\fancyhf{}
\cfoot{\thepage}
\setlength{\headheight}{16pt}

%%%%%%%%%%%%%%%%%%%%%%%%%%%%%%%%%%%%%%%%%%%%%%%%%%%%%
%%  Definition of the cover page and title page  %%
%%%%%%%%%%%%%%%%%%%%%%%%%%%%%%%%%%%%%%%%%%%%%%%%%%%%%
 
\ifthenelse{\boolean{englishLanguage}}{
%%% ENGLISH VERSION
\newcommand{\JMUTitle}[9]{

  \thispagestyle{empty}
  \vspace*{\stretch{1}}
  {\parindent0cm
  \rule{\linewidth}{.7ex}
  }
  \begin{flushright}
    \sffamily\bfseries\Huge
    #1\\
    \vspace*{\stretch{1}}
    \sffamily\bfseries\large
    #2\\
    \vspace*{\stretch{1}}
    \sffamily\bfseries\small
    #3
  \end{flushright}
  \rule{\linewidth}{.7ex}

  \vspace*{\stretch{1}}
  \begin{center}
    \includegraphics[width=2in]{seal.pdf} \\
    \vspace*{\stretch{1}}
    \Large #5 \\

    \vspace*{\stretch{2}}
   \large Chair for Enterprise Artificial Intelligence\\
    \large Universität Würzburg\\
    \vspace*{\stretch{1}}
    \large Supervisor:  #8 \\[1mm]
    \large Assistant:  #9 \\[1mm]
    \vspace*{\stretch{1}}
    \large #6, #7
  \end{center}
}}{
%%% GERMAN VERSION
\newcommand{\JMUTitle}[9]{

  \thispagestyle{empty}
  \vspace*{\stretch{1}}
  {\parindent0cm
  \rule{\linewidth}{.7ex}
  }
  \begin{flushright}
    \sffamily\bfseries\Huge
    #1\\
    \vspace*{\stretch{1}}
    \sffamily\bfseries\large
    #2\\
    \vspace*{\stretch{1}}
    \sffamily\bfseries\small
    #3
  \end{flushright}
  \rule{\linewidth}{.7ex}

  \vspace*{\stretch{1}}
  \begin{center}
    \includegraphics[width=2in]{seal.pdf} \\
    \vspace*{\stretch{1}}
    \Large #5 \\

    \vspace*{\stretch{2}}
   \large Lehrstuhl f\"ur K\"unstliche Intelligenz im Unternehemen\\
    \large Universität Würzburg\\
    \vspace*{\stretch{1}}
    \large Betreuer:  #8 \\[1mm]
    \large Assistent:  #9 \\[1mm]
    \vspace*{\stretch{1}}
    \large #6, den #7
  \end{center}
}} 


\titlespacing*{\section}
{0pt}{3.5ex plus 1ex minus .2ex}{.2ex plus .2ex}
\titlespacing*{\subsection}
{0pt}{1.5ex plus 1ex minus .2ex}{.2ex plus .2ex}
\titlespacing*{\subsubsection}
{0pt}{1.5ex plus 1ex minus .2ex}{.2ex plus .2ex}

%%%%%%%%%%%%%%%%%%%%%%%%%%%%
%%  Beginning of the document  %%
%%%%%%%%%%%%%%%%%%%%%%%%%%%%


\begin{document}
    \ifthenelse{\boolean{englishLanguage}}{
    \JMUTitle
      {TITLE }        % Title of the paper
      {FIRST NAME LAST NAME }                        % First and last name of the author
      {MATRICLE NUMBER}
      
      {Seminar Paper / Bachelor Thesis / Master Thesis } % Type of the work
      {W\"urzburg}                           % Place
      {01.01.2024}                          % Date of Submission
      {Prof. Dr. Gunther Gust}           % Name of the first examiner
      {Ignacio Úbeda/ Justus Ameling/ ...} % Name of the supervisor
    }{\JMUTitle
      {TITEL }        % Title of the paper
      {VORNAME NACHNAME }                        % First and last name of the author
      {MATRIKELNUMMER}
      
      {Seminararbeit / Bachelorarbeit / Masterarbeit} % Type of the work
      {W\"urzburg}                           % Place
      {01.01.2024}                          % Date of Submission
      {Prof. Dr. Gunther Gust}           % Name of the first examiner
      {Ignacio Úbeda/ Justus Ameling/ ...} % Name of the supervisor
      }
      
  \clearpage

\lhead{}
\pagenumbering{Roman} 
    \setcounter{page}{1}

\tableofcontents
\clearpage

\addcontentsline{toc}{section}{\listfigurename}
\listoffigures

\addcontentsline{toc}{section}{\listtablename}
\listoftables
\clearpage

\setlength{\parskip}{0.5em} 


%%%%%%%%%%%%%%%%%%%%%%%%%%%%
%%  Short Summary   %%
%%%%%%%%%%%%%%%%%%%%%%%%%%%%
\lhead{Abstract}
\section*{Abstract}

A brief summary of the procedure and the main results.

General Characteristics 
\begin{itemize}
    \item Objectivity: It should not contain any personal evaluation.
    \item Brevity: It should be as short as possible.
    \item Comprehensibility: It has a clear, comprehensible language and structure.
    \item Completeness: All essential facts should be included.
    \item Precision: It should accurately reflect the content and opinion of the original work.
\end{itemize}{}


%%%%%%%%%%%%%%%%%%%%%%%%%%%%
%%  Settings  %%
%%%%%%%%%%%%%%%%%%%%%%%%%%%%
\clearpage
\pagenumbering{arabic}  
    \setcounter{page}{1}
\lhead{\nouppercase{\leftmark}}

%%%%%%%%%%%%%%%%%%%%%%%%%%%%
%%  Main Section  %%
%%%%%%%%%%%%%%%%%%%%%%%%%%%%

\section{Section} \label{Introduction}
A main section - ideally, no sections should be empty.

\subsection{Subsection}
A subsection - ideally, no sections should be empty.

\subsubsection{Sub-Subsection}
A sub-subsection - ideally, no sections should be empty.

\section{Simple Style Sheets}

\textbf{This is a bold text}.

\textit{This is text in italics}.


Lists are often helpful for structuring:
\begin{itemize}
    \item First Entry
    \item Second Entry 
\end{itemize}

Numbered lists are often helpful for sequences:
\begin{enumerate}
    \item First Entry
    \item Second Entry 
\end{enumerate}


\section{Citation and referencing}

Articles in Journals like \citet{clemen1989combining} or conference papers like \citet{he2017mask} are cited as shown in this example. In other cases, the author might want to cite inside of brackets \citep{clemen1989combining}, even with multiple authors \citep{clemen1989combining,baumol1958warehouse,he2017mask}.

A page number must be given for monographs \citep[S. 28]{chollet2018deep}.

This is how a web source is cited: \citet{shiny1}. However, it may also be sufficient to provide the address in footnotes\footnote{\url{https://shiny.rstudio.com/tutorial/written-tutorial/lesson1/}} for brief information on the Internet.

This is how other parts of the work are referenced: Chapter \ref{Introduction}, equation \ref{eq:1} shows...



\section{Figures}

Figures require the package \textit{graphicx}. 
Ideally, vector graphics or high-resolution bitmaps should be used. 
A good option is to use PDFs.

\begin{figure}[h]
    \centering
    \includegraphics[width=0.3\textwidth]{seal.pdf}
    \caption{Seal of the university}
    \label{fig:my_label}
\end{figure}


\section{Tables}

The tabular environment specifies the number of columns, their width and any intermediate lines. 


\begin{table}[ht]
    \centering
    \caption{My Table}
        \begin{tabular}{ cccc } 
        \toprule
        col1 & col2 & col3 \\
        \midrule
        \multirow{3}{4em}{Multiple row} & cell2 & cell3 \\ 
        & cell5 & cell6 \\ 
        & cell8 & cell9 \\ 
        \bottomrule
    \end{tabular}
    \label{tab:countries}
\end{table}

\section{Equations}

\begin{equation}
    \sum_{i=1}^N x_i
    \label{eq:1}
\end{equation}



%%%%%%%%%%%%%%%%%%%%%%%%%%%%
%% The bibliography is inserted automatically
%%%%%%%%%%%%%%%%%%%%%%%%%%%%
\clearpage
\lhead{}
\printbibliography
\addcontentsline{toc}{section}{\bibname}


%%%%%%%%%%%%%%%%%%%%%%%%%%%%
%% Appendix
%%%%%%%%%%%%%%%%%%%%%%%%%%%%
\clearpage
\appendix
\section{Appendix A}

%%%%%%%%%%%%%%%%%%%%%%%%%%%%
%% Affidavit must be adapted 
%% in declaration.tex
%%%%%%%%%%%%%%%%%%%%%%%%%%%%
\input{declaration.tex}

\end{document}


\end{document}